
\documentclass[11pt,a4paper]{article}

\usepackage[margin=1in]{geometry}
\usepackage{amsmath,amssymb}
\usepackage{booktabs}
\usepackage{array}
\usepackage{enumitem}
\usepackage{microtype}
\usepackage[hidelinks]{hyperref}
\usepackage{setspace}

\setstretch{1.08}
\setlength{\parindent}{0pt}
\setlength{\parskip}{0.6em}

\title{\textbf{Progress Report: Development of a Computational Framework for Selective Radioprotection Using Natural Products}}
\author{
Edward Gyasi\\
Ghana Atomic Energy Commission\\
Radiation Protection Institute
}
\date{\today}

\begin{document}

\maketitle

\section{Project Overview}

This project aims to develop an integrated computational framework for prioritizing natural products with potential \textit{selective radioprotective} properties. The central scientific problem is to identify compounds that may support protection or compatibility of normal tissue during irradiation while retaining, or at least not compromising, inhibitory effects against cancer cells.

The project was motivated by experimental observations involving \textit{Moringa oleifera} extracts, in which differential effects were observed between normal fibroblast and cancer-cell systems following irradiation. Rather than assuming that \textit{Moringa} itself is radioprotective, the present work seeks to determine which individual phytochemicals or related natural products could plausibly contribute to such selective biological behaviour.

The long-term objective is therefore:

\begin{quote}
\textit{To develop and validate a computational framework that integrates cancer-cell inhibition, normal-cell compatibility, radiation-response biology and predictive confidence to identify natural-product candidates with potential selective radioprotective activity, using \textit{Moringa oleifera} phytochemicals as a biologically motivated case study.}
\end{quote}

\section{Scientific Strategy}

The project is being developed in sequential components so that each stage answers a distinct biological question before integration.

\begin{enumerate}[leftmargin=1.6em]
    \item \textbf{Cancer-potency modelling:} Can molecular structure predict inhibitory activity against MCF-7 breast cancer and DU145 prostate cancer cells?
    \item \textbf{Normal-cell compatibility:} Which compounds are compatible with normal fibroblast survival and show low normal-cell toxicity?
    \item \textbf{Radiation-response context:} Which candidate mechanisms are compatible with biological responses observed after irradiation in normal and cancer cells?
    \item \textbf{Integration:} Can these independent signals be combined into a selective-radioprotection ranking?
    \item \textbf{Natural-product screening:} Which natural products rank most favourably when the framework is applied prospectively?
    \item \textbf{\textit{Moringa oleifera} evaluation:} Are \textit{Moringa} phytochemicals individually or collectively enriched among the highest-ranked candidates?
\end{enumerate}

\section{Data Used to Date}

The current development stage focuses on the cancer-potency component.

\subsection{NCI-60 / CellMiner}

NCI-60 CellMiner compound-response data were used as the primary source for cancer-cell activity. Two cell lines were selected because they correspond to the cancer models of interest in the broader study:

\begin{itemize}[leftmargin=1.6em]
    \item \textbf{MCF-7}: human breast cancer cell line;
    \item \textbf{DU145}: human prostate cancer cell line.
\end{itemize}

The principal supervised endpoint was defined as absolute activity expressed as

\[
-\log_{10}(GI_{50}),
\]

where larger values indicate greater cancer-cell sensitivity to the compound.

Processed CellMiner activity z-scores were retained as secondary information describing relative cell-line sensitivity within the wider NCI-60 response profile. These z-scores were not used as substitutes for absolute potency.

\subsection{Chemical Structure Information}

Chemical identities were standardized using canonical molecular representations. Exact chemical duplicates were resolved before modelling so that the same chemical structure could not appear independently in different data partitions.

Molecular features currently include:

\begin{itemize}[leftmargin=1.6em]
    \item Morgan circular fingerprints (radius 2, 2048 bits, chirality enabled);
    \item a compact physicochemical descriptor set including molecular weight, cLogP, polar surface area, hydrogen-bond features, rotatable bonds, ring count, heavy-atom count and fraction Csp$^3$.
\end{itemize}

\section{Methodological Development Completed}

\subsection{Data Acquisition and Audit}

The NCI-60 raw and processed CellMiner datasets were acquired and audited before model development. Raw activity profiles were investigated to determine whether repeated rows represented genuine independent experiments.

The audit showed that all multi-row NSC records had identical full NCI-60 activity profiles. The repeated rows were therefore treated as duplicated profile representations rather than independent biological replicates.

A reconstruction analysis further demonstrated strong agreement between z-scores calculated from the raw NCI-60 activity profiles and the official processed CellMiner values. This supported the use of absolute $-\log_{10}(GI_{50})$ activity as the primary cancer-potency endpoint.

\subsection{Chemical Identity Resolution}

The modelling cohort was grouped by exact chemical identity. Duplicate chemical groups were audited for response agreement.

Primary duplicate discordance was defined only from the absolute cancer-potency endpoints. Exact-structure groups with a response range greater than 1.0 log$_{10}$ GI$_{50}$ unit in either MCF-7 or DU145 were excluded from the primary supervised cohort. Secondary CellMiner z-score discordance was retained only as a quality-control flag and did not determine primary-cohort membership.

\subsection{Final Frozen Cancer Cohort}

The resulting frozen primary cancer dataset contains:

\begin{center}
\begin{tabular}{lr}
\toprule
\textbf{Component} & \textbf{Count} \\
\midrule
Unique chemicals & 20,047 \\
Training set & 16,038 \\
Validation set & 2,004 \\
Test set & 2,005 \\
\bottomrule
\end{tabular}
\end{center}

The cohort was partitioned by Bemis--Murcko scaffold using a deterministic scaffold-group allocation procedure. The split was designed to preserve approximately 80/10/10 proportions for both molecules and scaffold groups.

No exact chemical identity or scaffold was shared between training, validation and test partitions.

\subsection{Baseline Machine-Learning Models}

The current benchmark evaluates conventional regression models before introducing more complex neural or graph-based approaches.

The models are:

\begin{itemize}[leftmargin=1.6em]
    \item DummyRegressor;
    \item Ridge regression;
    \item Elastic Net;
    \item Random Forest;
    \item Extra Trees;
    \item XGBoost.
\end{itemize}

Each method is being evaluated independently for MCF-7 and DU145 using both physicochemical descriptors and Morgan fingerprints.

Hyperparameters are selected using the validation set only, with mean absolute error (MAE) defined in advance as the primary model-selection criterion. The scaffold-separated test set remains untouched during model selection.

Performance metrics include:

\begin{itemize}[leftmargin=1.6em]
    \item MAE;
    \item RMSE;
    \item $R^2$;
    \item Pearson correlation;
    \item Spearman rank correlation;
    \item prediction bias;
    \item top-5\% and top-10\% overlap;
    \item enrichment factors at the top 5\% and 10\%.
\end{itemize}

The ranking metrics are especially important because the final application is prospective virtual screening rather than exact prediction of every GI$_{50}$ value.

\section{Progress and Preliminary Results}

The cancer-potency benchmark is still running; therefore, results presented here are preliminary and should not be interpreted as the final model-selection outcome.

For MCF-7, the null model produced a validation MAE of approximately 0.93 log units. Linear models based on simple physicochemical descriptors produced only modest improvements.

Nonlinear descriptor-based models performed better, with Random Forest and Extra Trees reducing validation MAE to approximately 0.81.

A stronger improvement was observed when molecular structure was represented using Morgan fingerprints. Preliminary results for the Morgan-fingerprint Random Forest include:

\begin{center}
\begin{tabular}{lr}
\toprule
\textbf{Validation metric} & \textbf{Preliminary value} \\
\midrule
MAE & 0.721 \\
RMSE & 0.939 \\
$R^2$ & 0.374 \\
Spearman $\rho$ & 0.622 \\
Top-5\% enrichment factor & 9.63 \\
Top-5\% overlap & 48.5\% \\
\bottomrule
\end{tabular}
\end{center}

These preliminary results suggest that molecular structure contains meaningful information about MCF-7 cancer-cell potency under scaffold-separated evaluation.

The results also indicate that Morgan fingerprints are more informative than broad physicochemical descriptors alone, suggesting that specific molecular substructures and their combinations are important determinants of activity.

The high top-5\% enrichment is particularly relevant to the long-term screening objective. Even if exact quantitative potency prediction remains moderate, strong enrichment means that the model may still be useful for prioritizing a relatively small fraction of candidate compounds for downstream investigation.

The complete MCF-7 benchmark, the full DU145 benchmark, final test-set results and sensitivity analyses remain under development.

\section{Current Achievements}

The principal achievements to date are:

\begin{enumerate}[leftmargin=1.6em]
    \item establishment of a defensible absolute cancer-potency endpoint using CellMiner NCI-60 activity data;
    \item validation of the raw-profile structure and clarification that repeated raw rows are not independent replicates;
    \item development of a curated MCF-7/DU145 chemical-response dataset;
    \item exact chemical identity resolution and exclusion of severely discordant duplicate labels;
    \item construction of molecular fingerprints and physicochemical descriptors;
    \item creation of a leakage-free scaffold-based train/validation/test split;
    \item implementation of a reproducible conventional machine-learning benchmark;
    \item preliminary evidence that structure-based models can meaningfully prioritize potent MCF-7 compounds under unseen-scaffold conditions.
\end{enumerate}

\section{Roadmap Toward the Final Selective-Radioprotection Framework}

The remaining work is organized as a connected sequence in which each stage contributes one necessary component of the final selective-radioprotection decision.

\subsection{Phase 1: Complete and Freeze the Cancer-Potency Component}

The immediate task is to complete the conventional MCF-7 and DU145 benchmarks and identify the best-performing structure-based representation and algorithm using the predefined validation procedure.

The resulting cancer component will provide a compound-level score:

\[
C_i = f_{\mathrm{cancer}}(x_i),
\]

where $C_i$ represents predicted anticancer potency for compound $i$ and $x_i$ represents its molecular structure.

If conventional fingerprints leave substantial unexplained signal, learned molecular representations such as graph neural networks may then be evaluated. A GNN will therefore be introduced only if it addresses a demonstrated limitation of the conventional baseline.

\subsection{Phase 2: Develop the Normal-Cell Compatibility Component}

The next stage will curate public compound-response data from suitable normal fibroblast or related normal-cell systems.

The aim will be to estimate a normal-cell compatibility term,

\[
N_i = f_{\mathrm{normal}}(x_i),
\]

representing low toxicity or favourable compatibility with normal-cell survival.

Where direct compound-response measurements in irradiated fibroblasts are unavailable, this stage will be interpreted specifically as normal-cell safety or compatibility rather than direct proof of radioprotection.

\subsection{Phase 3: Add Radiation-Response Biological Context}

Radiation-specific transcriptomic datasets will then be used to characterize biological pathways altered after irradiation in:

\begin{itemize}[leftmargin=1.6em]
    \item normal fibroblasts;
    \item MCF-7 cells;
    \item DU145 cells.
\end{itemize}

These expression datasets will not be treated as drug-response labels. Instead, they will define radiation-response signatures involving processes such as DNA-damage response, oxidative stress, apoptosis, cell-cycle control, repair and inflammatory signalling.

Candidate compounds will subsequently be assessed for mechanistic compatibility with these radiation-response signatures, producing a radiation-context component:

\[
R_i = f_{\mathrm{radiation}}(m_i,S_{\mathrm{rad}}),
\]

where $m_i$ represents compound-associated mechanistic information and $S_{\mathrm{rad}}$ represents the relevant radiation-response signature.

\subsection{Phase 4: Integrate the Biological Components}

Once the cancer, normal-cell and radiation-response components have been independently developed, they will be combined into an integrated selective-radioprotection framework.

Conceptually, a final score may take the form:

\[
\mathrm{SRS}_i =
w_C C_i +
w_N N_i +
w_R R_i -
w_T T_i -
w_U U_i,
\]

where:

\begin{itemize}[leftmargin=1.6em]
    \item $C_i$ = cancer-inhibition evidence;
    \item $N_i$ = normal-cell compatibility;
    \item $R_i$ = radiation-response compatibility;
    \item $T_i$ = toxicity penalty;
    \item $U_i$ = uncertainty or applicability-domain penalty.
\end{itemize}

The weighting scheme will be biologically justified and assessed through sensitivity analysis. It will not be optimized to force known \textit{Moringa} compounds toward the top of the ranking.

\subsection{Phase 5: Screen Natural-Product Chemical Space}

The completed framework will then be applied prospectively to a broad natural-product library, principally COCONUT.

Each natural product will be evaluated for:

\[
\text{cancer activity}
+
\text{normal-cell compatibility}
+
\text{radiation compatibility}
+
\text{prediction confidence}.
\]

Applicability-domain assessment will be included so that compounds lying far outside the training chemical space are appropriately penalized or flagged.

\subsection{Phase 6: Evaluate \textit{Moringa oleifera} as a Biological Case Study}

Reported phytochemicals from \textit{Moringa oleifera} seed and leaf material will be curated, structurally resolved and processed using the same final framework applied to the wider natural-product library.

Two questions will then be addressed:

\begin{enumerate}[leftmargin=1.6em]
    \item Which individual \textit{Moringa} phytochemicals receive the strongest selective-radioprotection predictions?
    \item Are \textit{Moringa} phytochemicals statistically enriched among the highest-ranked natural products?
\end{enumerate}

Enrichment will be examined at ranking thresholds such as the top 1\%, 5\% and 10\%, with permutation-based or related statistical testing used to determine whether any observed enrichment exceeds random expectation.

This stage will connect the computational framework back to the biological observation that originally motivated the project.

\subsection{Final Stage: Prioritization for Experimental Validation}

The final computational output will not simply be a ranked list of thousands of molecules. The aim is to nominate a small set of structurally diverse, high-confidence candidates for experimental follow-up.

Candidate selection will consider:

\begin{itemize}[leftmargin=1.6em]
    \item predicted cancer inhibition;
    \item normal-cell compatibility;
    \item radiation-response compatibility;
    \item model confidence and applicability domain;
    \item structural diversity;
    \item natural-product provenance;
    \item experimental feasibility.
\end{itemize}

A final shortlist of approximately 10--20 compounds would provide a practical basis for targeted experimental validation.

\section{Overall Development Pathway}

The full research programme can therefore be summarized as:

\[
\boxed{
\begin{array}{c}
\text{Cancer potency}\\
\downarrow\\
\text{Normal-cell compatibility}\\
\downarrow\\
\text{Radiation-response biology}\\
\downarrow\\
\text{Integrated selective-radioprotection score}\\
\downarrow\\
\text{Natural-product screening}\\
\downarrow\\
\text{\textit{Moringa oleifera} enrichment analysis}\\
\downarrow\\
\text{Experimental candidate prioritization}
\end{array}
}
\]

\section{Expected Scientific Contribution}

The expected contribution is a generalizable computational strategy for natural-product prioritization in selective radioprotection.

Rather than assuming that a compound is beneficial because it protects normal cells or originates from a medicinal plant, the framework will require simultaneous evidence that a candidate:

\begin{itemize}[leftmargin=1.6em]
    \item retains favourable anticancer characteristics;
    \item is compatible with normal-cell survival;
    \item has a biologically plausible relationship with radiation-response pathways;
    \item lies within an acceptable prediction domain;
    \item remains high-ranking when compared with a broad natural-product background.
\end{itemize}

This makes \textit{Moringa oleifera} a testable biological case study within a broader discovery framework rather than a predetermined positive outcome.

\section{Conclusion}

The project remains under active development. The cancer-potency component has progressed from raw data acquisition and endpoint validation to a frozen, leakage-free molecular modelling cohort and an active conventional machine-learning benchmark.

Preliminary MCF-7 results indicate meaningful structure--activity signal and strong enrichment of potent compounds among top model predictions. Completion of the MCF-7 and DU145 benchmarks will establish the first validated component of the broader selective-radioprotection framework.

The subsequent phases will extend the work from cancer potency to normal-cell compatibility, radiation-response biology, integrated natural-product screening and finally \textit{Moringa oleifera}-specific enrichment and experimental candidate prioritization.

\end{document}
